\chapter{Many-sorted discretization transformations}
\label{ch:ms_discrete}

Every many-sorted $n$-category can be embedded into a higher-dimensional $m$-category by extending
the family of sets with copies of the top-level set: dimensions above $n$ are filled with the
$n$-morphism set, and the extra operations are declared trivial. This construction extends naturally
to the $\omega$-categorical setting. Moreover, any many-sorted $n$-functor automatically preserves
this extended structure, giving rise to a functor between the resulting higher-dimensional
categories, with the extra dimensions acted upon by the top-level component of the original functor.
We call this process the \emph{discretization} transformation, and we show that it is functorial and
that discretization transformations at different dimensions commute with one another.

\section{Discrete many-sorted categories}

We begin by defining the extended family of sets and then construct a many-sorted category structure
on it.

\begin{definition}[Discrete many-sorted finite family]
  \label{def:ManySortedDiscreteFamily}
  \uses{def:ManySortedNCategory}
  Let $\cat{C}$ be a many-sorted $n$-category on a family $(C_k)_{k \leq n}$, and let $m \in \NN$
  with $n < m$. The \emph{discrete finite family} $(C^{>m}_k)_{k \leq m}$ is defined by
  \[
    C^{>m}_k =
    \begin{cases}
      C_k & \text{if } k \leq n, \\
      C_n & \text{if } n < k \leq m.
    \end{cases}
  \]
\end{definition}

\begin{construction}[Discrete many-sorted finite category]
  \label{def:ManySortedNCategory_discrete}
  \uses{def:ManySortedDiscreteFamily, def:ManySortedNCategory}
  \notready
  Let $\cat{C}$ be a many-sorted $n$-category on a family $(C_k)_{k \leq n}$, and let $m \in \NN$
  with $n < m$. The \emph{discrete $m$-category} of $\cat{C}$ is the many-sorted $m$-category
  $\cat{C}^{>m}$ on the discrete family $(C^{>m}_k)_{k \leq m}$ of
  Definition~\ref{def:ManySortedDiscreteFamily}, where the operations at each pair $(k, j)$ with $j
  < k$ are defined as follows:

  \begin{itemize}
    \item If $k \leq n$ (and hence $j < k \leq n$): the source $\Sc^{(k,j)}$, target $\Tg^{(k,j)}$,
      identity $\Idm^{(k,j)}$, and composition $\circ^{(k,j)}$ are those of $\cat{C}$.
    \item If $k > n$ and $j \leq n$: we have $C^{>m}_k = C_n$ and $C^{>m}_j = C_j$, so the source
      and target maps $\Sc^{(k,j)}, \Tg^{(k,j)} \colon C_n \to C_j$ are defined as $\Sc^{(n,j)}$ and
      $\Tg^{(n,j)}$ respectively; the identity map $\Idm^{(k,j)} \colon C_j \to C_n$ is defined as
      $\Idm^{(n,j)}$; and composition $\circ^{(k,j)}$ is defined as $\circ^{(n,j)}$.
    \item If $j > n$ (and hence $k > j > n$): we have $C^{>m}_k = C^{>m}_j = C_n$. The source and
      target maps are the identity on $C_n$, the identity map $\Idm^{(k,j)}$ is also the identity,
      and composition is defined only between equal morphisms and returns the morphism itself.
  \end{itemize}
\end{construction}

\begin{proof}
  Each axiom, whether a pre-category axiom at a pair $(k, j)$ or a cross-dimensional axiom at a
  triple $i < j < k$, either involves only indices at most $n$, in which case it reduces to the
  corresponding axiom of $\cat{C}$; involves indices above $n$ whose operations coincide with those
  at the top dimension $n$, in which case it again follows from an axiom of $\cat{C}$; or involves
  only indices above $n$, where all operations are trivial and the axiom holds immediately.
\end{proof}

\begin{definition}[Discrete many-sorted omega family]
  \label{def:ManySortedDiscreteOmegaFamily}
  \uses{def:ManySortedNCategory}
  Let $\cat{C}$ be a many-sorted $n$-category on a family $(C_k)_{k \leq n}$. The \emph{discrete
  omega family} $(C^{>\omega}_k)_{k \in \NN}$ is defined by
  \[
    C^{>\omega}_k =
    \begin{cases}
      C_k & \text{if } k \leq n, \\
      C_n & \text{if } k > n.
    \end{cases}
  \]
\end{definition}

\begin{construction}[Discrete many-sorted omega category]
  \label{def:ManySortedNCategory_discreteOmega}
  \uses{def:ManySortedDiscreteOmegaFamily, def:ManySortedNCategory, def:ManySortedOmegaCategory}
  \notready
  Let $\cat{C}$ be a many-sorted $n$-category on a family $(C_k)_{k \leq n}$. The \emph{discrete
  $\omega$-category} of $\cat{C}$ is the many-sorted $\omega$-category $\cat{C}^{>\omega}$ on the
  discrete omega family $(C^{>\omega}_k)_{k \in \NN}$ of
  Definition~\ref{def:ManySortedDiscreteOmegaFamily}, defined by the same case distinction as in
  Definition~\ref{def:ManySortedNCategory_discrete}, now with the indices $k$ and $j$ ranging over
  all of $\NN$.
\end{construction}

\begin{proof}
  The many-sorted category axioms follow by the same argument as in
  Definition~\ref{def:ManySortedNCategory_discrete}.
\end{proof}

\section{Discrete many-sorted functors}

\begin{construction}[Discrete many-sorted finite functor]
  \label{def:ManySortedNFunctor_discrete}
  \uses{def:ManySortedNFunctor, def:ManySortedNCategory_discrete}
  \notready
  Let $\cat{C}$ and $\cat{D}$ be many-sorted $n$-categories, and let $F = (F_k)_{k \leq n} \colon
  \cat{C} \to \cat{D}$ be a many-sorted $n$-functor. For any $m \in \NN$ with $n < m$, the family
  $(F^{>m}_k)_{k \leq m}$ defined by
  \[
    F^{>m}_k =
    \begin{cases}
      F_k & \text{if } k \leq n, \\
      F_n & \text{if } n < k \leq m
    \end{cases}
  \]
  is a many-sorted $m$-functor $F^{>m} \colon \cat{C}^{>m} \to \cat{D}^{>m}$.
\end{construction}

\begin{proof}
  Each preservation axiom at a pair $(k, j)$ with $j < k$ either involves only indices at most $n$,
  where it follows from the functor axioms of $F$; involves indices above $n$ whose maps and
  operations coincide with those at the top dimension $n$, where it again follows from a functor
  axiom of $F$; or involves only indices above $n$, where all maps coincide and all operations are
  trivial, so the axiom holds immediately.
\end{proof}

\begin{construction}[Discrete many-sorted omega functor]
  \label{def:ManySortedNFunctor_discreteOmega}
  \uses{def:ManySortedOmegaFunctor, def:ManySortedNCategory_discreteOmega}
  \notready
  Let $\cat{C}$ and $\cat{D}$ be many-sorted $n$-categories, and let $F = (F_k)_{k \leq n} \colon
  \cat{C} \to \cat{D}$ be a many-sorted $n$-functor. The family $(F^{>\omega}_k)_{k \in \NN}$
  defined by
  \[
    F^{>\omega}_k =
    \begin{cases}
      F_k & \text{if } k \leq n, \\
      F_n & \text{if } k > n
    \end{cases}
  \]
  is a many-sorted $\omega$-functor $F^{>\omega} \colon \cat{C}^{>\omega} \to \cat{D}^{>\omega}$.
\end{construction}

\begin{proof}
  The argument is identical to that of Definition~\ref{def:ManySortedNFunctor_discrete}.
\end{proof}

\section{Many-sorted discretization functors}

The discrete constructions of the previous sections are in fact functorial, as the following
definition makes precise.

\begin{construction}[Many-sorted discretization functor]
  \label{def:ManySortedDiscretizationFunctor}
  \uses{
    def:ManySortedNCategory_discrete,
    def:ManySortedNCategory_discreteOmega,
    def:ManySortedNFunctor_discrete,
    def:ManySortedNFunctor_discreteOmega
  }
  \notready
  Let $n, m \in \ENat$ with $n \leq m$. The \emph{many-sorted discretization functor}
  $\DiscreteMS^{(n, m)}$ from $n\CatMS$ to $m\CatMS$ is the functor that maps every many-sorted
  $n$-category $\cat{C}$ to its discrete $m$-category $\cat{C}^{>m}$ and every many-sorted
  $n$-functor $F \colon \cat{C} \to \cat{D}$ to its discrete $m$-functor $F^{>m} \colon \cat{C}^{>m}
  \to \cat{D}^{>m}$. If $n = m$, then $\DiscreteMS^{(n, n)}$ stands for the identity functor on
  $n\CatMS$.
\end{construction}

\begin{proof}
  Functoriality follows from the fact that the discrete constructions preserve the underlying maps:
  the identity functor on $\cat{C}$ yields the identity functor on $\cat{C}^{>m}$, and the
  composition of two $n$-functors yields the composition of their discrete $m$-functors.
\end{proof}

Moreover, discretization transformations at different dimensions are compatible with one another.

\begin{theorem}[Compatibility of many-sorted discretization functors]
  \label{thm:ManySortedDiscretizationFunctor_compatibility}
  \uses{def:ManySortedDiscretizationFunctor}
  \notready
  The assignment $(n, m) \mapsto \DiscreteMS^{(n, m)}$ defines a functor from $(\ENat, \leq)$ to
  $\cat{Cat}$. That is, for all $k \leq n \leq m$ in $\ENat$,
  \[
    \DiscreteMS^{(n, m)} \circ \DiscreteMS^{(k, n)} =
    \DiscreteMS^{(k, m)}.
  \]
\end{theorem}

\begin{proof}
  Both sides produce the same family of sets and the same operations. The categorical structures
  coincide because extending first to dimension $n$ and then to $m$ introduces the same trivial
  structure on dimensions above $k$ as extending directly to $m$.
\end{proof}
